Online review platforms are a primary source of information for consumers choosing
where to eat, yet the same restaurant is often rated differently on each of them.
This project builds an integrated database of restaurants in Milan from three
platforms --- Google Maps, Tripadvisor, and TheFork --- and uses it to study how
consistent restaurant ratings and metadata are across sources.

A complete ELT pipeline was developed. Data are acquired through the Google Places
API and two dedicated web scrapers, loaded as raw documents into MongoDB, cleaned
and normalized per source, and then integrated. Integration relies on schema
matching, a Google-anchored entity-resolution stage with calibrated thresholds,
dedicated handling of chain and franchise venues, and an optional large-language-model
adjudication of the uncertain cases. The resulting Google-seeded collection preserves
each platform's ratings side by side and is projected into a flat ClickHouse layer for
analytical querying. The quality of the source data is assessed before integration,
while the entity-resolution process that drives the integration is evaluated
afterwards against a manually labelled gold standard.

The final dataset contains 10,054 restaurants. Ratings are generally consistent
across platforms, but Tripadvisor is systematically more severe than Google and
TheFork; rating level and cross-platform agreement vary by cuisine, with mainstream
categories more contested than specialist ones; and data completeness is strongly
spatially structured, decreasing from the historic centre towards the periphery.
